\documentclass[10pt,a4paper]{amsart}
\usepackage{amsmath,amssymb,amsthm,graphicx}
\usepackage{verbatim}
\newenvironment{reference}{\comment}{\endcomment}
%\newenvironment{reference}{}{}
\newenvironment{slogan}{\comment}{\endcomment}
\newenvironment{history}{\comment}{\endcomment}

% For commutative diagrams we use Xy-pic
\usepackage[all]{xy}

% We use 2cell for 2-commutative diagrams.
\xyoption{2cell}
% Theorem environments.
%
\theoremstyle{plain}
\newtheorem{theorem}[subsection]{Theorem}
\newtheorem{proposition}[subsection]{Proposition}
\newtheorem{lemma}[subsection]{Lemma}

\theoremstyle{definition}
\newtheorem{definition}[subsection]{Definition}
\newtheorem{example}[subsection]{Example}
\newtheorem{exercise}[subsection]{Exercise}
\newtheorem{situation}[subsection]{Situation}

\theoremstyle{remark}
\newtheorem{remark}[subsection]{Remark}
\newtheorem{remarks}[subsection]{Remarks}

\numberwithin{equation}{subsection}

% Macros
%
\def\lim{\mathop{\mathrm{lim}}\nolimits}
\def\colim{\mathop{\mathrm{colim}}\nolimits}
\def\Spec{\mathop{\mathrm{Spec}}}
\def\Hom{\mathop{\mathrm{Hom}}\nolimits}
\def\Ext{\mathop{\mathrm{Ext}}\nolimits}
\def\SheafHom{\mathop{\mathcal{H}\!\mathit{om}}\nolimits}
\def\SheafExt{\mathop{\mathcal{E}\!\mathit{xt}}\nolimits}
\def\Sch{\mathit{Sch}}
\def\Mor{\mathop{\mathrm{Mor}}\nolimits}
\def\Ob{\mathop{\mathrm{Ob}}\nolimits}
\def\Sh{\mathop{\mathit{Sh}}\nolimits}
\def\NL{\mathop{N\!L}\nolimits}
\def\CH{\mathop{\mathrm{CH}}\nolimits}
\def\proetale{{pro\text{-}\acute{e}tale}}
\def\etale{{\acute{e}tale}}
\def\QCoh{\mathit{QCoh}}
\def\Ker{\mathop{\mathrm{Ker}}}
\def\Im{\mathop{\mathrm{Im}}}
\def\Coker{\mathop{\mathrm{Coker}}}
\def\Coim{\mathop{\mathrm{Coim}}}

% Boxtimes
%
\DeclareMathSymbol{\boxtimes}{\mathbin}{AMSa}{"02}

%
% Macros for moduli stacks/spaces
%
\def\QCohstack{\mathcal{QC}\!\mathit{oh}}
\def\Cohstack{\mathcal{C}\!\mathit{oh}}
\def\Spacesstack{\mathcal{S}\!\mathit{paces}}
\def\Quotfunctor{\mathrm{Quot}}
\def\Hilbfunctor{\mathrm{Hilb}}
\def\Curvesstack{\mathcal{C}\!\mathit{urves}}
\def\Polarizedstack{\mathcal{P}\!\mathit{olarized}}
\def\Complexesstack{\mathcal{C}\!\mathit{omplexes}}
% \Pic is the operator that assigns to X its picard group, usage \Pic(X)
% \Picardstack_{X/B} denotes the Picard stack of X over B
% \Picardfunctor_{X/B} denotes the Picard functor of X over B
\def\Pic{\mathop{\mathrm{Pic}}\nolimits}
\def\Picardstack{\mathcal{P}\!\mathit{ic}}
\def\Picardfunctor{\mathrm{Pic}}
\def\Deformationcategory{\mathcal{D}\!\mathit{ef}}
\usepackage{cleveref}
\crefname{originaldefinition}{original definition}{original definitions}
\Crefname{originaldefinition}{Original definition}{Original definitions}
\makeatletter
\newlabel{definition-weaker-than-canonical}{{00WQ}{source}{}{sourcepointer-1}{}}
\newlabel{definition-weaker-than-canonical@cref}{{[originaldefinition][0][]00WQ}{[1][0][]source}{}{}{}}
\makeatother
\crefname{originallemma}{original lemma}{original lemmas}
\Crefname{originallemma}{Original lemma}{Original lemmas}
\makeatletter
\newlabel{lemma-describe-j-shriek}{{03CD}{source}{}{sourcepointer-2}{}}
\newlabel{lemma-describe-j-shriek@cref}{{[originallemma][0][]03CD}{[1][0][]source}{}{}{}}
\makeatother
\title{No sheafification is needed for the lower shriek over a subcanonical site}
\author{The Stacks Project Authors}
\date{stacks-project:sites.tex:local-snapshot-20261012; GNU FDL 1.2 or later}
\begin{document}
\maketitle
\noindent\textit{Editorial note.} All delivered mathematical text is the author's own, copied from the pinned Stacks Project chapter file sites-cohomology.tex; no statement or proof step was rewritten, reordered or removed.\par\medskip
\noindent\textit{Editorial note.} The proof is the author's complete argument; every step the author writes out is inside the counted span and no closing sentence is a placeholder.\par\medskip
\noindent\textit{Editorial note.} Every labelled reference occurring inside the counted spans is carried in assembly.external\_references as a non-load-bearing navigation pointer (same-chapter ones at their definition site, cross-chapter ones at the invocation inside the counted span), so no load-bearing mathematics is replaced by a stub. No bibliography citation occurs inside the counted spans.\par\medskip
\noindent\textit{Editorial note.} The checked chapter list for same-name collisions is the set of stacks chapters staged in this repository: chow, cohomology, divisors, etale-cohomology, flat, more-algebra, more-morphisms, morphisms, spaces-more-morphisms, varieties together with this card's chapter sites-cohomology.\par\medskip
\noindent\textit{Editorial note.} The preamble is the chapter's own retained minimal preamble (sources/stacks-sites-cohomology/retained-preamble.tex), a verbatim excerpt of the native preamble: the Xy-pic setup and the native theorem declarations with the mathematical macros.\par\medskip
\noindent\textit{Allowed standard prerequisites (editorial).}\begin{itemize}
\item The description of the lower shriek along a localization. Same chapter, Lemma lemma-describe-j-shriek. Supplies the presheaf whose sheafification is the lower shriek.
\item Weakly canonical topologies. Same chapter, Definition definition-weaker-than-canonical. Supplies the condition under which the gluing of the maps is unique.
\end{itemize}\medskip
\ifdefined\hypertarget\hypertarget{sourcepointer-1}{}\fi
\noindent\textit{Source reference.} Original definition 00WQ: sources/stacks-sites/pinned-original-source.tex, source lines 6272--6272. Navigation pointer invoked at proof line 6272; the counted span contains the pointer only and the cited result is not reproduced inside it, so the pointer is registered as a non-load-bearing navigation stub in the established form for Stacks material. No bibliography citation is involved.\par\medskip
\ifdefined\hypertarget\hypertarget{sourcepointer-2}{}\fi
\noindent\textit{Source reference.} Original lemma 03CD: sources/stacks-sites/pinned-original-source.tex, source lines 6281--6281. Navigation pointer invoked at proof line 6281; the counted span contains the pointer only and the cited result is not reproduced inside it, so the pointer is registered as a non-load-bearing navigation stub in the established form for Stacks material. No bibliography citation is involved.\par\medskip
\begin{lemma}
\label{lemma-describe-j-shriek-good-site}
Let $\mathcal{C}$ be a site.
Let $U \in \Ob(\mathcal{C})$.
If the topology on $\mathcal{C}$ is subcanonical, see
Definition \ref{definition-weaker-than-canonical},
and if $\mathcal{G}$ is a sheaf on $\mathcal{C}/U$, then
$$
j_{U!}(\mathcal{G})(V)
=
\coprod\nolimits_{\varphi \in \Mor_\mathcal{C}(V, U)}
\mathcal{G}(V \xrightarrow{\varphi} U),
$$
in other words sheafification is not necessary in
Lemma \ref{lemma-describe-j-shriek}.
\end{lemma}

\begin{proof}
Let $\mathcal{V} = \{V_i \to V\}_{i \in I}$ be a covering of $V$
in the site $\mathcal{C}$.
We are going to check the sheaf condition for the presheaf $\mathcal{H}$
of Lemma \ref{lemma-describe-j-shriek} directly.
Let $(s_i, \varphi_i)_{i \in I} \in \prod_i \mathcal{H}(V_i)$,
This means $\varphi_i : V_i \to U$ is a morphism in $\mathcal{C}$, and
$s_i \in \mathcal{G}(V_i \xrightarrow{\varphi_i} U)$.
The restriction of the pair $(s_i, \varphi_i)$ to
$V_i \times_V V_j$ is the pair
$(s_i|_{V_i \times_V V_j/U}, \text{pr}_1 \circ \varphi_i)$, and
likewise the restriction of the pair $(s_j, \varphi_j)$ to
$V_i \times_V V_j$ is the pair
$(s_j|_{V_i \times_V V_j/U}, \text{pr}_2 \circ \varphi_j)$.
Hence, if the family $(s_i, \varphi_i)$ lies in
$\check{H}^0(\mathcal{V}, \mathcal{H})$, then we see that
$\text{pr}_1 \circ \varphi_i = \text{pr}_2 \circ \varphi_j$.
The condition that the topology on $\mathcal{C}$ is weaker than the canonical
topology then implies that there exists a unique morphism
$\varphi : V \to U$ such that $\varphi_i$ is the composition
of $V_i \to V$ with $\varphi$. At this point the sheaf condition for
$\mathcal{G}$ guarantees that the sections $s_i$ glue to a unique
section $s \in \mathcal{G}(V \xrightarrow{\varphi} U)$.
Hence $(s, \varphi) \in \mathcal{H}(V)$ as desired.
\end{proof}

\end{document}
